% ═══════════════════════════════════════════════════════════════════════
\section{Contractions}
% ═══════════════════════════════════════════════════════════════════════

\begin{frame}{Définition 2.14 --- Contraction}
  \small Les exemples précédents invitent à regarder les fonctions qui
  rapprochent les points.
  \begin{Enonce}{Définition 2.14 --- Contraction}
    Soit $(E,d)$ un espace métrique et $f:E\to E$. $f$ est une
    \defi{contraction} sur $E$ s'il existe $k\in\intfo{0}{1}$ tel que
    \[
      d\big(f(x),f(y)\big)\leq k\cdot d(x,y)\qquad\forall\,x,y\in E.
    \]
    La constante $k$ est le \defi{rapport de contraction} de $f$.
  \end{Enonce}
  \small Chaque application de $f$ multiplie les distances par au plus
  $k<1$.
\end{frame}

\begin{frame}{Exemple 2.10 --- Contractions du plan}
  \small
  \begin{itemize}
    \item Une homothétie de rapport $k$, $|k|<1$ : contraction de
      rapport $|k|$ sur $(\R^2,\dE)$.
    \item Plus généralement, toute similitude de rapport $k<1$ (comme
      celle de l'exemple 2.9).
    \item Mais une contraction n'est pas forcément une
      similitude :
      \[
        f\col{x}{y} = \mat{1/2}{0}{0}{2/3}\col{x}{y}
      \]
      est une contraction de rapport $2/3$ (elle réduit davantage
      selon $x$ que selon $y$), sans préserver les angles.
  \end{itemize}
  \FigBox[0.28\textheight]{%
    \begin{tikzpicture}[scale=1.1]
      \begin{scope}
        \fill[fig] \motifF;
        \node[etiq] at (0.4,-0.3) {motif};
      \end{scope}
      \draw[fleche] (1.2,0.5) -- node[etiq, above]{$f$} (2.0,0.5);
      \begin{scope}[xshift=2.5cm, xscale=0.5, yscale=0.6667]
        \fill[figim] \motifF;
      \end{scope}
      \node[etiq] at (2.7,-0.3) {image};
    \end{tikzpicture}}
\end{frame}

\begin{frame}{Remarque 2.10 --- Le rôle de la borne stricte $k<1$}
  \begin{Rem}{Remarque 2.10}
    La condition $k<1$, et non $k\leq1$, est essentielle.
    \begin{itemize}
      \item $f(x)=x+1$ sur $(\R,\dE)$ : $d(f(x),f(y)) = d(x,y)$. Elle
        conserve les distances ($k=1$)… et n'a aucun point fixe
        (exemple 2.8 (i)).
      \item $f(x)=3x+2$ : $d(f(x),f(y)) = 3\,d(x,y)$. Non seulement
        $k=3\geq1$, mais $f$ éloigne les points d'un facteur $3$.
    \end{itemize}
    \vspace{2pt}
    C'est ce qui explique que l'itération de l'exemple 2.8 (iii)
    s'éloigne de son unique point fixe $x=-1$ au lieu d'y converger.
  \end{Rem}
\end{frame}
